\section{Conclusion}
A useful energy measurement requires more than a sufficiently high sampling rate. Across six Jetson workloads, fixed 2 kS/s meets the empirical 1\% energy criterion at the tested 2-, 5- and 10-s intervals, while a 100-ms interval requires a persistent 80-kS/s minimum. The separate FP16 study shows why long energy integration and fluctuation analysis need markedly different rates. The Hailo extension shows why the next decision is about the deployed observation: electrical boundary, actual interval and workload determine what apparent meter agreement means.

Embedded sensing agrees closely with external DC acquisition during the tested long continuous workloads, while vendor and AC comparisons remain specific to their domains and transformations. Duration-wise Hailo pairing makes the distinction tangible: short-request mean-power agreement can conceal a much larger energy difference when the retained spans differ. Spectral characterisation and fixed-rate protocol controls add complementary evidence that neither a stable average nor a faster meter alone establishes representative execution or temporal fidelity.

The practical recommendation is therefore to validate the measurement configuration as a whole: define the quantity and boundary, test the rate against its interval, compare energy together with power and span, and verify the execution conditions. This supports repeatable, workload-specific energy assessment without turning one successful operating point into a universal sampling rule or meter-accuracy ranking.
